\subsection{Sequential Compactness}

\defn{Subsequence}{
    Let $x_n$ be a sequence in some $X$.

    If $n_1 < n_2 < n_3 < \dots$ is a sequence of strictly increasing indices, then the sequence $x_{n_1},x_{n_2},\dots$ is a subsequence of $x_n$.
}

\defn{Sequentially Compact}{
    A topological space $X$ is sequentially compact if every sequence of $X$ has a convergent subsequence.
}

\thmp{}{
    Suppose $X$ is a metrizable space. Then $X$ is compact if and only if $X$ is sequentially compact.
}{
    $(\implies)$

    Suppose $X$ is not sequentially compact, then let $x_n$ be a sequence with no convergent subsequence.

    Let $A$ be the set of elements appearing in the sequence $x_n$.

    \clmp{}{$A$ is closed}{
        Let $x \in \bar{A}$. Then, because $X$ is metrizable, $x$ is a limit of a sequence of elements $y_n \in A$.

        \noindent If $x \notin A$, then a subsequence of $y_n$ is also a subsequence of $x_n$.

        \noindent Know $y_1=x_{n_1}$. Assume,
        \[y_{m_1}=x_{n_1},\dots,y_{m_j}=x_{n_j}\]
        for indices
        \[m_1<m_2,\dots<m_j\]
        \[n_1<n_2<\dots<n_j\]
        The set $\{y_m \mid m > m_j\}$ is infinite since if it were finite we would get a constant subsequence which implies that $x \in A$

        \noindent This intersects $A \setminus \{x_1,\dots,x_{n_j}\}$.

        \noindent So take $y_{m_{j+1}}=x_{n_{j+1}}$ to be any elt in this set.

        \noindent Build a subseq. of $y_n$ converging with subsequence of $x_n$ This is a contradiction thus $x \in A \implies \bar{A}=A$ thus $A$ is closed.
    }
    For each $m$ we know $x_m$ is not the limit of a subsequence of $x_n$. So, there exists a neighborhood of $x_m$ which contains only finitely many elements from $A$.

    Also know $X$ is Hausdorff, so there is a neighborhood $U_n$ of $x_n$ s.t. $U_n$ contains no other elements of $A$.

    Covering:
    \[(X \setminus A) \sqcup \bigcup_{n=1}^\infty U_n\] has no finite subcover. Thus, $X$ is not compact.

    $(\impliedby)$
    Assume $X$ is sequentially compact and metrizable.
    \clmp{1}{
        Any open covering of $X$ has a Lebesgue number.
    }{
        Let $\mathcal{C}$ be an open cover of $X$. Assume for contradiction no Lebesgue number exists.

        \noindent So, for all $n \in \N$ there is a set $D_n$ s.t. diam$(D_n) < 1/n$, and $D_n \nsubseteq C$ for any $C \in \mathcal{C}$.

        \noindent Fix $x_n \in D_n$. Let $x$ be the limit of a subsequence of $x_n$. $x \in C \subseteq \mathcal{C}$, so $x \in B(x,\epsilon) \subseteq C$ for some $\epsilon > 0$.

        \noindent For some large $n$, we have $d(x,x_n)< \epsilon/2$ and $1/n < \epsilon/2$.

        \noindent $x_n \in D_n$, and diam$(D_n) < 1/n < \epsilon/2$ by triangle inequality, $D_n \subseteq B(x,\epsilon) \subseteq C$. This is a contradiction.
    }
    \clmp{2}{
        Given any $\epsilon > 0$, there exists a finite collection of points $x_1,\dots,x_n$ s.t.
        \[\bigcup_{i=1}^nB(x_i,\epsilon)=X\]
    }{
        Assume for contradiction that the claim is not possible. Fix $x_1 \in X$. $B(x_1, \epsilon) \neq X$, so let $x_2 \in X \setminus B(x_1,\epsilon)$.

        \noindent Now define sequence s.t. $x_i \notin B(x_1,\epsilon) \cup \dots \cup B(x_{i-1},\epsilon)$.
        
        \noindent We know that the union of these $j$ balls is not equal to $X$, so let $x_{j+1}$ be in the complement of the union of the $j$ balls.

        \noindent This sequence $x_n$ has no convergent subsequence since $d(x_i,x_j) > \epsilon$ for all $i \neq j$. This means that no subsequence is Cauchy. Thus, no subsequence converges.
    }

    Let $\mathcal{C}$ be an open covering of $X$. Let $\delta > 0$ be a Lebesgue number for the covering. Cover $X$ by finitely many balls with radius $\epsilon=\delta/3$.

    Thus, diam$(B(x,\epsilon)) \leq 2\epsilon = 2\delta/3 < \delta$, so each ball in this finite covering is contained in some $C \in \mathcal{C}$.

    Pick some $C$ containing each ball in the cover by $\epsilon$-balls, so this is a finite subscover of $X$ by sets in $\mathcal{C}$.
}

\thm{Well-Ordering theorem}{
    Every set can be equipped with a well-ordering
}

\cor{
    There exists an uncountable well-ordered set.
}

\thmp{}{
    There exists a well-ordered uncountable set $\omega_1$ s.t. every interval $(-\infty, a) \subset \omega_1$ is countable $[0,a)$
}{
    Let $S$ be any uncountable well-ordered set
    \clmp{}{
        There exists another uncountable well-ordered set $T$ which contains an uncountable interval
    }{
        Let $T=S \times \{1,2\}$ with dictionary order.
    }
    The interval $[0,z)$ is uncountable if $z=(b,2)$ for $b \in S$. 

    Let $\Omega$ smallest element $z \in T$ s.t. $[0,z)$ is uncountable.

    Define $\omega_1=[0,\Omega), \overline{\omega_1}=[0,\Omega]$

    Every interval in $[0,a)$ in $\omega_1$ is countable.
}

\propp{
    Every countable subset of $\omega_1$ is bounded.
}{
    Let $A \subset \omega_1$ be countable. Let $B=\displaystyle\bigcup_{a \in A}[0,a]$.

    $B$ is countable, so $B \subsetneq \omega_1$. $\exists z \in \omega_1 \setminus B$.

    $B \subseteq [0,z)$ if $\exists y \in B$, $y \geq z$, then $[0,y] \subset B$, but $y \in [0,y]$ and $z \notin B$.
    
    Thus, $B$ is upper bounded since $\overline{\omega_1}$ is well-ordered, so it satisfies least upper bound property. Thus, it is closed and bounded, so it is compact.

    $\omega_1$ is not compact.
}